\chapter{Introdução}

Assistentes pessoais por voz tornaram-se parte do cotidiano em telefones celulares, caixas de som
e automóveis. Em computadores pessoais, porém, sobretudo nos que utilizam o sistema operacional
Linux, a oferta de assistentes integrados ao ambiente de trabalho ainda é limitada. As soluções
comerciais costumam depender de serviços em nuvem permanentemente à escuta, oferecem pouca
personalização da aparência e da voz e não guardam um histórico verificável das atividades
realizadas pelo usuário. Ao mesmo tempo, o surgimento de modelos de linguagem de grande porte
(\en{Large Language Models}, LLM), capazes de interpretar pedidos em linguagem natural e de acionar
funções externas, reduziu drasticamente o esforço necessário para construir um assistente que
compreenda frases livres, e não apenas comandos rígidos \cite{vaswani2017}.

Construir um assistente de voz útil, entretanto, envolve muito mais do que conectar um reconhecedor
de fala a um modelo de linguagem. É preciso capturar o áudio com qualidade, decidir quando o usuário
terminou de falar, sintetizar a resposta sem longas esperas, sincronizar a boca de um avatar com a
voz, persistir dados de forma confiável e manter a interface gráfica fluida enquanto diversas tarefas
concorrentes acontecem. Esse conjunto de preocupações, distribuído entre várias tecnologias, torna o
software suscetível a acoplamento excessivo e à degradação de sua estrutura à medida que novas
funções são acrescentadas, fenômeno descrito por \textcite{martin2019} como o principal fator do custo
crescente de manutenção.

Este trabalho descreve a concepção, a construção incremental e a validação do \textbf{Robofael}, um
assistente pessoal de atividades diárias que roda localmente em um computador com Linux. O Robofael
é representado por um avatar animado criado a partir de fotografias do próprio usuário, ouve e
responde por voz, organiza tarefas e lembretes, responde perguntas sobre o que foi feito em
determinado período, abre programas instalados e conversa em português ou em inglês. Além da
descrição funcional, o trabalho dá ênfase à qualidade interna do software: a solução foi organizada
segundo a arquitetura hexagonal \cite{cockburn2005} e, ao final, revisada de ponta a ponta à luz dos
princípios SOLID \cite{martin2003} e das práticas de código limpo \cite{martin2009}, com regras de
projeto verificadas automaticamente por testes.

\section{Problema de pesquisa}

O problema que orienta este trabalho pode ser enunciado da seguinte forma: \emph{como construir um
assistente pessoal por voz para computadores com Linux que seja responsivo, preserve a privacidade do
usuário, mantenha um histórico confiável de suas atividades e, ao mesmo tempo, conserve uma estrutura
interna fácil de entender, testar e evoluir?}

\section{Objetivos}

\subsection{Objetivo geral}

Desenvolver e validar um assistente pessoal de atividades diárias, com avatar animado e interação por
voz, cuja arquitetura siga os princípios SOLID e as práticas de código limpo, com comportamento e
regras de projeto verificados por testes automatizados.

\subsection{Objetivos específicos}

\begin{enumerate}[label=\alph*)]
  \item gerar, a partir de fotografias do usuário, um avatar em formato de busto com movimentação da
        boca sincronizada à voz sintetizada;
  \item implementar o ciclo de conversa por voz (ouvir, compreender e responder) com latência
        percebida próxima à da primeira frase da resposta;
  \item permitir a ativação sem as mãos por uma frase de chamada reconhecida localmente;
  \item oferecer gestão de tarefas e lembretes e uma memória de longo prazo auditável, capaz de
        responder o que foi feito em um dia, em uma semana ou em outro período;
  \item integrar o assistente ao ambiente Linux, permitindo abrir programas instalados por comando de
        voz;
  \item suportar os idiomas português do Brasil e inglês dos Estados Unidos, com troca em tempo de
        execução;
  \item disponibilizar um instalador padrão para o Ubuntu;
  \item revisar a solução segundo os princípios SOLID e as práticas de código limpo, tornando as
        regras de arquitetura verificáveis por testes automatizados;
  \item documentar os requisitos funcionais e não funcionais e todo o processo de construção.
\end{enumerate}

\section{Justificativa}

A relevância do trabalho manifesta-se em três dimensões. Na dimensão do usuário, um assistente que
roda no próprio computador, escuta passivamente apenas de forma local e guarda os dados em um arquivo
sob controle do usuário atende a uma preocupação crescente com privacidade. A memória de longo prazo
auditável, por sua vez, transforma o assistente em uma ferramenta de acompanhamento de produtividade:
ao perguntar “o que eu fiz na semana passada?”, o usuário obtém uma resposta baseada em registros
cuja integridade pode ser verificada.

Na dimensão técnica, o projeto reúne, em um único sistema de porte moderado, problemas clássicos de
engenharia de software: concorrência entre captura de áudio, geração de texto, síntese de voz e
animação; integração com serviços externos sujeitos a falhas; persistência transacional; e
internacionalização. Isso o torna um estudo de caso adequado para avaliar como princípios de projeto
orientado a objetos se comportam em um contexto real.

Na dimensão acadêmica, o trabalho documenta um processo de desenvolvimento conduzido com o apoio de
um agente de programação baseado em inteligência artificial e mostra que a qualidade estrutural, nesse
cenário, pode ser preservada por meio de verificações automáticas, como testes de arquitetura, limites
de complexidade e verificação estática de tipos, em vez de depender exclusivamente de disciplina
manual.

\section{Delimitação do escopo}

O sistema destina-se a um único usuário por conta do sistema operacional e foi construído e testado
no Ubuntu 24.04 com Python 3.12, embora o código seja compatível com Python 3.10 ou superior. O
reconhecimento de fala para comandos usa o serviço de reconhecimento do Google acessado por meio da
biblioteca SpeechRecognition, e a síntese de voz utiliza vozes neurais do serviço edge-tts, com o gTTS
como reserva; ambos exigem conexão com a internet. A interpretação em linguagem natural usa a API do
Claude ou, sem ela, um modelo local servido pelo Ollama, mais lento quando executado no processador; o
modo offline cobre um conjunto definido de comandos. O treinamento de
modelos acústicos próprios, a clonagem completa do timbre de voz e o suporte a outros sistemas
operacionais estão fora do escopo deste trabalho.

\section{Organização do trabalho}

Além desta introdução, o trabalho está organizado em nove capítulos. O Capítulo 2 apresenta a
fundamentação teórica sobre assistentes de voz, processamento de fala, modelos de linguagem,
arquitetura de software, princípios SOLID e código limpo. O Capítulo 3 descreve a metodologia de
desenvolvimento adotada. O Capítulo 4 especifica os requisitos funcionais e não funcionais. O
Capítulo 5 detalha a arquitetura e os fluxos de execução. O Capítulo 6 relata, passo a passo, a
construção incremental do sistema. O Capítulo 7 apresenta a revisão do código à luz dos princípios
SOLID e do código limpo. O Capítulo 8 trata da estratégia de testes e da validação. O Capítulo 9
discute os resultados obtidos, e o Capítulo 10 apresenta as conclusões e os trabalhos futuros.
